\begin{problem}[第36届全国中学生物理竞赛决赛][介质半圆柱体中的电场与电路]{112}
    一个长为 $l$、内外半径分别为 $a$ 和 $b$ 的半圆柱体由两种不同的耗损电介质构成，它们的相对介电常数和电导率分别为：$\varepsilon_{r1}$ 和 $\sigma_{1}$（$0\leq\varphi<\theta_{0}$ 区域），$\varepsilon_{r2}$ 和 $\sigma_{2}$（$\theta_{0}<\varphi\leq\pi$ 区域），其横截面如图 a 所示；在半圆柱两侧底部镀有金属膜，两金属膜间加有直流电压 $V_{0}$，并达到稳定。已知真空介电常量为 $\varepsilon_{0}$，两种介质的相对介电常数大、电导率小。忽略边缘效应。
    \begin{subproblem}
        求介质内电场强度和电势的分布；
    \end{subproblem}
    \begin{subproblem}
        求 $\varphi = \theta_0$ 界面处的总电荷；
    \end{subproblem}
    \begin{subproblem}
        分别计算 $0 \leq \varphi < \theta_{0}$ 和 $\theta_{0} < \varphi \leq \pi$ 两介质区域的电阻和电容；
    \end{subproblem}
    \begin{subproblem}
        在 $t = 0$ 时刻断开电源，求随后的两金属膜间电压随时间的变化（该过程可视为似稳过程），并画出相应的等效电路图。
    \end{subproblem}
\end{problem}